\section{Related Work}

\para{Inference-time Steering of Robot Policies}
Recent work on inference-time adaptation for robot policies typically follow three paradigms: sample-then-verify~\cite{wu2025foresight,kwok2025robomonkey,hansteer2026,yuan2026actasklearnuncertaintyaware}, noise-space optimization~\cite{wagenmaker2025steering}, and classifier-based guidance~\cite{sun2026latent,du2025dynaguidesteeringdiffusionpolices,wang2024inference}. While effective, they all have focused on visual guidance signals that provide global task context but often miss fine-grained contact details in manipulation~\cite{zhang2026touchguide,xue2025reactive,higuera2026visuo,cui2026vitacman}. Some tactile steering works instead use touch to provide direct contact information~\cite{tian2019manipulation,zhang2026touchguide}, but do not capture global progress. Specifically, TouchGuide~\cite{zhang2026touchguide} learns action-space verifiers via contrastive learning, but it requires costly demonstrations to obtain positive samples and does not explicitly model dynamics, limiting its scalability.
Instead, we propose a visuo-tactile steering approach that operates at different temporal scales with a learned latent world model. In our work, visual lookahead provides global long-horizon steering, while tactile refinement enables short-horizon, contact-sensitive corrections.




\para{Tactile Sensing for Robot Decision Making}
We use vision-based tactile sensors, specifically GelSight Mini~\cite{yuan2017gelsight} due to its dense measurements of contact geometry, deformation, and force-related signals. The existing datasets with GelSight~\cite{yuan2017gelsight,yuan2015measurement,helmut2025learning} have been used to pretrain encoders for tactile images~\cite{fenganytouch,fenganytouch2,higuerasparsh,higuera2025tactile,xie2026universal}, providing semantically meaningful latent spaces for tactile reasoning. Tactile sensing has also been widely used for contact-rich manipulation~\cite{pmlr-v78-calandra17a,calandra2018more,lee2019making,chen2023visuo, zhu2025touch}, and recent work incorporates tactile feedback into learned policies for fine-grained reactive control~\cite{huang20253d,heng2025vitacformerlearningcrossmodalrepresentation,xue2025reactive,chen2026multimodal}, primarily during policy training or low-level control. In contrast, we use tactile as an inference-time guidance signal: candidate actions are scored by aligning predicted tactile embeddings with textual contact goals, enabling scalable tactile verification without demonstrations.





\para{World Modeling in Robot Decision Making}
Dynamics models have been studied in planning, control, and policy learning for robotics~\cite{sutton1991dyna,deisenroth2011pilco,chua2018deep}. Recent visual world models focus on learning latent dynamics ~\cite{zhou2025dino,huang2026vjepa,assran2025v}, adapting diffusion or flow-based video models~\cite{guo2025ctrl,alhaija2025cosmos,gao2026dreamdojo} or training with model-based reinforcement learning~\cite{hafner2019dream,hafner2023mastering,wu2023daydreamer, hafner2025training}. Meanwhile, recent visuo-tactile world models incorporate touch into dynamics modeling~\cite{higuera2026visuo,tian2019manipulation,zheng2026omnivta,yuan2026vtam}.
In this work, we train a visuo-tactile world model to evaluate the multimodal outcomes of actions directly in the latent space. 
Therefore, actions can be scored directly in the latent space without decoding images, enabling efficient text-conditioned 
guidance.
















